\chapter{Policy gradients and actor-critic methods}
\label{ch:pg}

Value-based methods learn $q_\star$ and act greedily. Policy-based methods
take the other road: they parametrise the policy itself, $\pi_\theta(a \mid
s)$, and improve it by gradient ascent on the expected return. Combined
with a learnt value function (the \emph{critic}) that reduces the variance
of the gradient estimates, this gives the actor-critic methods, A2C and its
successor PPO (\cref{ch:ppo}), which produced the best agent of the
repository.

\section{Parametrised policies}

The agents of the repository use a \emph{softmax policy}: a network outputs
one score (\emph{logit}) $z_\theta(s, a)$ per action, and
\begin{equation}
  \pi_\theta(a \mid s) = \frac{\exp z_\theta(s, a)}{\sum_{b}
  \exp z_\theta(s, b)}.
\end{equation}
Compared with an $\varepsilon$-greedy policy, a softmax policy changes
smoothly with $\theta$, can represent any stochastic policy, and explores
by itself in the states where it is uncertain.

\begin{proposition}{Score of a softmax policy}{softmax-score}
  The gradient of $\log \pi_\theta(a \mid s)$ with respect to the logits
  $z = (z(s, b))_b$ of state $s$ is $e_a - \pi_\theta(\cdot \mid s)$, where
  $e_a$ is the indicator vector of action $a$.
\end{proposition}

\begin{proof}
  $\log \pi_\theta(a \mid s) = z(s, a) - \log \sum_b \exp z(s, b)$, whose
  derivative with respect to $z(s, c)$ is $\mathbb{1}[c = a] - \exp z(s,
  c) / \sum_b \exp z(s, b) = \mathbb{1}[c = a] - \pi_\theta(c \mid s)$.
\end{proof}

Increasing $\log \pi_\theta(a \mid s)$ thus raises the logit of $a$ and
lowers the others in proportion to their probabilities.

\section{The policy gradient theorem}

We maximise the expected return from the initial state distribution
$\mu$,
\begin{equation}
  J(\theta) = \E_{\pi_\theta}\left[G_0\right]
  = \sum_s \mu(s)\, v_{\pi_\theta}(s).
\end{equation}
The difficulty is that $\theta$ influences the distribution of the
trajectories, through the actions, and the dynamics $p$ that also shape it
are unknown. The policy gradient theorem shows that the gradient
nevertheless has an expression that can be estimated from trajectories,
without $p$. We prove it for episodes of bounded length $H$, which covers
Pacman: a game is aborted after at most $1000$ moves.

\begin{theorem}{Policy gradient theorem}{pg}
  Suppose every episode ends (or is aborted) within $H$ steps, and $\pi_\theta$
  is differentiable in $\theta$. Then
  \begin{equation}
    \label{eq:pg}
    \nabla_\theta J(\theta) = \E_{\pi_\theta}\left[ \sum_{t=0}^{H-1}
    \nabla_\theta \log \pi_\theta(A_t \mid S_t)\, \gamma^t\, G_t \right]
    = \E_{\pi_\theta}\left[ \sum_{t=0}^{H-1} \nabla_\theta \log
    \pi_\theta(A_t \mid S_t)\, \gamma^t\, q_{\pi_\theta}(S_t, A_t) \right],
  \end{equation}
  where, after the end of an episode, the terms are $0$.
\end{theorem}

\begin{proof}
  \emph{Step 1 (likelihood ratio).} A trajectory $\tau = (s_0, a_0, r_1,
  \dots, s_H)$ has probability
  \begin{equation*}
    P_\theta(\tau) = \mu(s_0) \prod_{t=0}^{H-1} \pi_\theta(a_t \mid s_t)\,
    p(s_{t+1}, r_{t+1} \mid s_t, a_t),
  \end{equation*}
  and $G_0(\tau) = \sum_{t} \gamma^t r_{t+1}$ does not depend on $\theta$.
  There are finitely many trajectories, so we may differentiate under the
  sum: using $\nabla P = P \nabla \log P$ where $P > 0$,
  \begin{equation*}
    \nabla_\theta J = \sum_\tau G_0(\tau)\, \nabla_\theta P_\theta(\tau)
    = \E_{\pi_\theta}\left[ G_0\, \nabla_\theta \log P_\theta(\tau)
    \right]
    = \E_{\pi_\theta}\left[ G_0 \sum_{t=0}^{H-1} \nabla_\theta \log
    \pi_\theta(A_t \mid S_t) \right],
  \end{equation*}
  because $\mu$ and $p$ do not depend on $\theta$ and disappear from the
  gradient of the logarithm.

  \emph{Step 2 (causality).} Write $\psi_t = \nabla_\theta \log
  \pi_\theta(A_t \mid S_t)$. For $k \leq t$, the reward $R_k$ is determined
  by the history up to $S_t$, and given that history,
  \begin{equation*}
    \E\left[\psi_t \given S_0, \dots, S_t\right] = \sum_a \pi_\theta(a \mid
    S_t)\, \nabla_\theta \log \pi_\theta(a \mid S_t) = \nabla_\theta \sum_a
    \pi_\theta(a \mid S_t) = \nabla_\theta 1 = 0.
  \end{equation*}
  Hence $\E[R_k \psi_t] = 0$ for $k \leq t$: an action cannot influence the
  rewards received before it. Removing these terms from $G_0 \psi_t$ leaves
  $\sum_{k > t} \gamma^{k-1} R_k\, \psi_t = \gamma^t G_t\, \psi_t$, which
  proves the first equality.

  \emph{Step 3 (conditioning).} Given $(S_t, A_t)$, $\psi_t$ is fixed and,
  by the Markov property, $\E[G_t \mid S_0, A_0, \dots, S_t, A_t] =
  q_{\pi_\theta}(S_t, A_t)$. The tower property gives the second equality.
\end{proof}

\begin{remark}{The factor $\gamma^t$}{gamma-t}
  The factor $\gamma^t$ in \eqref{eq:pg} gives little weight to the late
  steps of an episode. Practical implementations, including the
  repository's, drop it and weight all the steps equally; they then follow
  a direction that is not exactly the gradient of $J$ (Thomas, 2014), but
  which treats every state as a possible starting point, which works well
  in practice.
\end{remark}

The simplest algorithm built on \eqref{eq:pg} is REINFORCE (Williams,
1992): play episodes and take the step $\theta \gets \theta + \alpha \sum_t
\psi_t G_t$. It is unbiased but very noisy, since each $G_t$ is a Monte
Carlo return.

\section{Baselines and advantages}

\begin{proposition}{Baselines}{baseline}
  For any function $b : \States \to \R$ (a \emph{baseline}),
  $\E_{\pi_\theta}\left[\nabla_\theta \log \pi_\theta(A_t \mid S_t)\,
  b(S_t)\right] = 0$. Consequently \eqref{eq:pg} remains true with $G_t$
  replaced by $G_t - b(S_t)$, or $q_{\pi_\theta}(S_t, A_t)$ by
  $q_{\pi_\theta}(S_t, A_t) - b(S_t)$.
\end{proposition}

\begin{proof}
  Conditionally on $S_t$, $b(S_t)$ is a constant and
  $\E[\psi_t \mid S_t] = 0$ (step 2 of the previous proof).
\end{proof}

A baseline does not change the expected gradient but can greatly reduce
its variance. The natural choice is $b = v_{\pi_\theta}$, which gives the
\emph{advantage}
\begin{equation}
  A_\pi(s, a) = q_\pi(s, a) - v_\pi(s),
\end{equation}
the amount by which action $a$ is better than the average action of the
policy in $s$. The gradient step then increases the probability of the
actions that did better than expected and decreases the others. The
advantage of an action is often much smaller than its value: in Pacman, the
value of a state may be $100$, while the difference between going east and
going west is a few points.

\section{Actor-critic methods}

The advantage is unknown; an \emph{actor-critic} method learns a value
function $V_w \approx v_{\pi_\theta}$ (the \emph{critic}) alongside the
policy (the \emph{actor}), and uses it to estimate the advantages.

\begin{proposition}{The TD error estimates the advantage}{td-advantage}
  Let $\delta_t = R_{t+1} + \gamma\, v_\pi(S_{t+1}) - v_\pi(S_t)$, with the
  true value function. Then $\E_\pi[\delta_t \mid S_t = s, A_t = a] =
  A_\pi(s, a)$.
\end{proposition}

\begin{proof}
  By \eqref{eq:q-from-v}, $\E_\pi[R_{t+1} + \gamma v_\pi(S_{t+1}) \mid S_t
  = s, A_t = a] = q_\pi(s, a)$; subtract $v_\pi(s)$.
\end{proof}

With the critic $V_w$ instead of $v_\pi$, $\delta_t$ is a biased but
low-variance estimate of the advantage. Monte Carlo advantages $G_t -
V_w(S_t)$ are unbiased (given a baseline) but noisy; the generalised
advantage estimator interpolates between the two.

\subsection{Generalised advantage estimation}

\begin{definition}{GAE}{gae}
  Let $\delta_t = R_{t+1} + \gamma V(S_{t+1}) - V(S_t)$ be the TD errors of
  a critic $V$ (with $V = 0$ after a terminal state). For $\lambda \in [0,
  1]$, the \emph{generalised advantage estimate} (Schulman et al., 2016) is
  \begin{equation}
    \label{eq:gae}
    \hat A_t^{\lambda} = \sum_{l=0}^{\infty} (\gamma \lambda)^l\,
    \delta_{t+l},
  \end{equation}
  where the sum stops at the end of the episode. It satisfies the recursion
  $\hat A_t^\lambda = \delta_t + \gamma \lambda\, \hat A_{t+1}^\lambda$, used
  to compute it backwards along a trajectory.
\end{definition}

\begin{proposition}{GAE interpolates between TD and Monte Carlo}{gae-props}
  For an episode ending at time $T$ (so that the sums are finite):
  \begin{enumerate}[nosep, label=(\roman*)]
    \item $\hat A_t^0 = \delta_t$;
    \item $\hat A_t^1 = G_t - V(S_t)$;
    \item $\hat A_t^\lambda + V(S_t) = G_t^\lambda$, the $\lambda$-return
      \eqref{eq:lambda-return}.
  \end{enumerate}
\end{proposition}

\begin{proof}
  (i) is immediate. (ii) The sum telescopes:
  \begin{equation*}
    \sum_{l=0}^{T-t-1} \gamma^l \left( R_{t+l+1} + \gamma V(S_{t+l+1}) -
    V(S_{t+l}) \right) = \sum_{l=0}^{T-t-1} \gamma^l R_{t+l+1} +
    \gamma^{T-t} V(S_T) - V(S_t) = G_t - V(S_t),
  \end{equation*}
  since $V(S_T) = 0$ at the terminal state. (iii) The same telescoping over
  the first $n$ terms gives $\sum_{l=0}^{n-1} \gamma^l \delta_{t+l} =
  G_{t:t+n} - V(S_t)$. Hence
  \begin{equation*}
    (1 - \lambda) \sum_{n \geq 1} \lambda^{n-1} \left( G_{t:t+n} - V(S_t)
    \right)
    = (1 - \lambda) \sum_{n \geq 1} \lambda^{n-1} \sum_{l=0}^{n-1}
      \gamma^l \delta_{t+l}
    = \sum_{l \geq 0} \gamma^l \delta_{t+l} (1 - \lambda) \sum_{n > l}
      \lambda^{n-1},
  \end{equation*}
  and $(1 - \lambda) \sum_{n > l} \lambda^{n-1} = \lambda^l$. The left-hand
  side is $G_t^\lambda - V(S_t)$ because the weights $(1 -
  \lambda)\lambda^{n-1}$ sum to $1$ (for $\lambda < 1$; the case
  $\lambda = 1$ is (ii)). All sums are finite, which justifies the
  exchange.
\end{proof}

The parameter $\lambda$ trades bias (small $\lambda$, which relies on the
critic) against variance (large $\lambda$, which relies on the sampled
rewards). The repository uses $\lambda = 0.95$, and the critic is trained
towards the targets $\hat A_t^\lambda + V(S_t)$, i.e.\ the
$\lambda$-returns.

\begin{lstlisting}[caption={GAE by the backward recursion
  (\code{algorithms.py}, \code{A2C.compute\_gae})}]
# rewards, values, dones: arrays [n_steps, n_envs]; dones[t] = 1 if the
# game ended after step t; last_values: value of the states reached
advantages = np.zeros_like(rewards)
last_gae = 0.
for t in reversed(range(n_steps)):
    next_values = last_values if t == n_steps - 1 else values[t + 1]
    not_done = 1. - dones[t]
    delta = rewards[t] + self.gamma * next_values * not_done - values[t]
    last_gae = delta + self.gamma * self.gae_lambda * not_done * last_gae
    advantages[t] = last_gae
return advantages, advantages + values    # value targets: lambda-returns
\end{lstlisting}
A game that is aborted (not ended) after step $t$ is marked as done, but
its reward is first increased by $\gamma V(S_{t+1})$: the bootstrap of
\cref{rem:truncation}.

\section{Continuous actions: Gaussian policies}
\label{sec:gaussian}

Policy gradients extend directly to continuous actions $a \in \R^k$, where
value-based methods would need a maximisation over a continuum. The
standard choice is a Gaussian policy whose mean is computed by a network,
with a diagonal covariance (state-independent, or also computed by the
network):
\begin{equation}
  \pi_\theta(a \mid s) = \mathcal{N}\left(a;\, \mu_\theta(s),\,
  \mathrm{diag}(\sigma_\theta(s)^2)\right).
\end{equation}

\begin{proposition}{Score of a Gaussian policy}{gaussian-score}
  In one dimension, with mean $\mu$ and standard deviation $\sigma =
  e^{\eta}$,
  \begin{equation}
    \frac{\partial \log \pi(a)}{\partial \mu} = \frac{a - \mu}{\sigma^2},
    \qquad
    \frac{\partial \log \pi(a)}{\partial \eta} = \frac{(a -
    \mu)^2}{\sigma^2} - 1.
  \end{equation}
\end{proposition}

\begin{proof}
  $\log \pi(a) = -\frac{(a - \mu)^2}{2 e^{2\eta}} - \eta - \frac{1}{2}
  \log(2\pi)$; differentiate with respect to $\mu$ and $\eta$.
\end{proof}

An action better than expected ($\hat A > 0$) moves the mean towards it;
the standard deviation grows if the action was far from the mean
($\abs{a - \mu} > \sigma$) and shrinks otherwise. Actions bounded in an
interval are usually obtained by squashing a Gaussian sample with $\tanh$,
with the corresponding correction of $\log \pi$ (\cref{ch:continuous}).

\section{The natural policy gradient}

The gradient $\nabla_\theta J$ is the direction of steepest ascent for the
Euclidean distance between \emph{parameters}, which has no meaning for the
\emph{policies}: rescaling a parameter changes the gradient, and a small
change of a parameter can change the policy a lot in some states and not at
all in others. The natural gradient measures steps by the change of the
policy itself, with the KL divergence.

\begin{definition}{Fisher information matrix}{fisher}
  For a policy $\pi_\theta$ and a state distribution $d$,
  \begin{equation}
    F(\theta) = \E_{s \sim d,\, a \sim \pi_\theta(\cdot \mid s)}\left[
    \nabla_\theta \log \pi_\theta(a \mid s)\, \nabla_\theta \log
    \pi_\theta(a \mid s)^\top \right].
  \end{equation}
\end{definition}

\begin{proposition}{KL divergence to second order}{kl-fisher}
  Let $\bar K(\theta') = \E_{s \sim d}[\mathrm{KL}(\pi_\theta(\cdot \mid
  s) \,\|\, \pi_{\theta'}(\cdot \mid s))]$. For smooth policies with
  finitely many actions, $\bar K(\theta + \Delta) = \frac{1}{2}
  \Delta^\top F(\theta) \Delta + O(\norm{\Delta}^3)$.
\end{proposition}

\begin{proof}
  Fix $s$ and write $\pi_\theta(a)$ for $\pi_\theta(a \mid s)$. The function
  $\theta' \mapsto \mathrm{KL}(\pi_\theta \| \pi_{\theta'}) = \sum_a
  \pi_\theta(a) (\log \pi_\theta(a) - \log \pi_{\theta'}(a))$ vanishes
  at $\theta' = \theta$. Its gradient there is $-\sum_a \pi_\theta(a)
  \nabla \log \pi_\theta(a) = -\nabla \sum_a \pi_\theta(a) = 0$. Its
  Hessian is $-\sum_a \pi_\theta(a) \nabla^2 \log \pi_\theta(a)$, and
  differentiating $\sum_a \pi_\theta(a) \nabla \log \pi_\theta(a) = 0$
  once more gives $\sum_a \pi_\theta(a) \left( \nabla^2 \log
  \pi_\theta(a) + \nabla \log \pi_\theta(a) \nabla \log
  \pi_\theta(a)^\top \right) = 0$: the Hessian is the Fisher matrix of
  state $s$. Average over $s$ and apply Taylor's theorem.
\end{proof}

\begin{proposition}{Natural gradient direction}{natural-gradient}
  If $F$ is positive definite, the maximiser of the linear approximation
  $\nabla J^\top \Delta$ under the quadratic constraint $\frac{1}{2}
  \Delta^\top F \Delta \leq \delta$ is
  \begin{equation}
    \Delta^\star = \sqrt{\frac{2 \delta}{g^\top F^{-1} g}}\, F^{-1} g,
    \qquad g = \nabla_\theta J.
  \end{equation}
\end{proposition}

\begin{proof}
  With $u = F^{1/2} \Delta$ the problem is to maximise $(F^{-1/2} g)^\top u$
  subject to $\norm{u}^2 \leq 2 \delta$. By Cauchy--Schwarz the maximiser is
  $u$ colinear with $F^{-1/2} g$ and of norm $\sqrt{2 \delta}$, i.e.\
  $\Delta = F^{-1/2} u \propto F^{-1} g$, with the stated scaling.
\end{proof}

The \emph{natural policy gradient} (Kakade, 2001) follows $F^{-1} \nabla
J$. It is invariant to reparametrisations of the policy, and the size of
its steps is measured in KL divergence, the quantity that controls how much
the state distribution can change (\cref{ch:ppo}). Computing $F^{-1} g$
exactly is impossible for large networks; TRPO approximates it with
conjugate gradients, and PPO replaces the constraint by clipping.

\section{A2C}

\emph{Advantage actor-critic} (A2C, the synchronous version of A3C, Mnih et
al., 2016) runs $N$ environments in parallel for $n$ steps, then performs
one update on the $N n$ collected transitions.

\begin{algobox}{A2C with GAE}
  \begin{algorithmic}[1]
    \For{each iteration}
      \State play $n$ steps in each of the $N$ environments with
        $\pi_\theta$; store $s_t, a_t, r_{t+1}, d_{t+1}$ and $V_w(s_t)$
      \State compute $\hat A_t$ by the backward recursion
        $\hat A_t = \delta_t + \gamma \lambda (1 - d_{t+1}) \hat A_{t+1}$,
        bootstrapping the last step with $V_w$
      \State normalise the advantages (mean $0$, standard deviation $1$)
      \State one gradient step on the loss
        \begin{equation*}
          \mathcal{L} = -\frac{1}{Nn} \sum_t \hat A_t \log \pi_\theta(a_t
          \mid s_t) + \frac{c_v}{2Nn} \sum_t \left( \hat A_t + V_{w,
          \text{old}}(s_t) - V_w(s_t) \right)^2 - c_e\, \bar{\mathcal{H}}
        \end{equation*}
    \EndFor
  \end{algorithmic}
\end{algobox}

The last term is an \emph{entropy bonus}: $\bar{\mathcal{H}}$ is the
average entropy $-\sum_a \pi_\theta(a \mid s_t) \log \pi_\theta(a \mid s_t)$
of the policy over the batch. It keeps the policy from becoming
deterministic too early, which would stop exploration. In the repository
the actor and the critic share the same network (two heads on a common
trunk), with $c_v = 0.5$ and $c_e = 0.03$.

\begin{pacman}{The entropy coefficient}{entropy}
  With $c_e = 0.01$, short test trainings of the repository's agent saw the
  entropy of the policy fall below $0.85$ (out of a maximum of $\log 5
  \approx 1.61$) while the agent still lost almost every game: the policy
  had become nearly deterministic before finding good behaviours. With
  $c_e = 0.03$ the entropy stabilised around $1.2$ at the same stage and the
  reward improved faster, which is why the repository uses $0.03$.
\end{pacman}

\section{The actor-critic network of the repository}

\begin{pacman}{Input, network, symmetry, curriculum}{a2c-repo}
  \begin{itemize}[nosep]
    \item \textbf{Input:} the feature planes of the last two images
      ($14 \times 9 \times 13$ in game \code{v1}).
    \item \textbf{Network} (\code{ActorCriticCNN}): the same full-resolution
      convolutional trunk as the DQN agent, a fully connected layer of $256$
      units, a policy head (five logits) and a value head.
    \item \textbf{Symmetry:} the maze is symmetric with respect to its
      vertical axis, and the network is made exactly equivariant to this
      symmetry (\cref{prop:equivariance}).
    \item \textbf{Curriculum:} the aggressiveness of the ghosts starts at
      $30\%$ of its value and increases linearly to its full value over the
      first $20\%$ of the training, so that the agent first survives long
      enough to discover the rewards.
  \end{itemize}
\end{pacman}

Let $M$ be the horizontal mirror of an input (reverse the order of the
columns of every plane) and $P$ the permutation of the actions that
exchanges \code{E} and \code{W} and fixes \code{N}, \code{S}, \code{I}.
Both are involutions: $M^2 = \mathrm{id}$ and $P^2 = \mathrm{id}$. A network
$f$ giving logits and a value, $f(s) = (z(s), V(s))$, is
\emph{equivariant} if $z(Ms) = P z(s)$ and $V(Ms) = V(s)$: the mirror image
of a situation gets the mirrored policy and the same value.

\begin{proposition}{Symmetrisation}{equivariance}
  For any network $f = (z, V)$, the symmetrised network
  \begin{equation}
    z_{\mathrm{sym}}(s) = \tfrac{1}{2} \left( z(s) + P z(Ms) \right), \qquad
    V_{\mathrm{sym}}(s) = \tfrac{1}{2} \left( V(s) + V(Ms) \right)
  \end{equation}
  is equivariant. If $f$ is already equivariant, $f_{\mathrm{sym}} = f$.
\end{proposition}

\begin{proof}
  Using $M^2 = \mathrm{id}$ and $P^2 = \mathrm{id}$,
  \begin{equation*}
    z_{\mathrm{sym}}(Ms) = \tfrac{1}{2} \left( z(Ms) + P z(s) \right)
    = \tfrac{1}{2} P \left( P z(Ms) + z(s) \right) = P z_{\mathrm{sym}}(s),
  \end{equation*}
  and $V_{\mathrm{sym}}(Ms) = \frac{1}{2}(V(Ms) + V(s)) =
  V_{\mathrm{sym}}(s)$. If $z(Ms) = P z(s)$ then $P z(Ms) = z(s)$ and
  $z_{\mathrm{sym}} = z$; similarly for $V$.
\end{proof}

The symmetrised network is evaluated on each state and its mirror image (in
a single batch), and trained as is. The agent can then never learn
different behaviours on the two halves of the maze, and every game teaches
it about both halves at once: this is equivalent to doubling the data. The
same construction works for any group of symmetries of the board, by
averaging over the group.

\begin{lstlisting}[caption={The symmetrised network (\code{neural\_net.py},
  \code{ActorCriticCNN.forward})}]
# action_mirror[i] = index of the mirror of action i (E <-> W)
logits, values = self._forward(
    torch.cat([state, torch.flip(state, dims=[-1])]))  # s and Ms
b = state.shape[0]
logits = (logits[:b] + logits[b:][:, self.action_mirror]) / 2  # z_sym
values = (values[:b] + values[b:]) / 2                         # V_sym
\end{lstlisting}

\section{Summary}

\begin{itemize}[nosep]
  \item The policy gradient theorem expresses $\nabla J$ as an expectation
    over trajectories that does not involve the dynamics.
  \item A baseline leaves the gradient unchanged and reduces its variance;
    with the value function it gives the advantage.
  \item Actor-critic methods learn the value function to estimate the
    advantages; GAE interpolates between the TD error and the Monte Carlo
    advantage with the parameter $\lambda$.
  \item A2C runs several environments in parallel and makes one update per
    batch of transitions; the entropy bonus preserves exploration.
\end{itemize}
A2C makes one gradient step per batch and throws the batch away: the next
chapter shows how PPO reuses each batch several times without letting the
policy change too much.
